% Generated draft: requires mathematical and source review.
\documentclass[11pt]{article}
\usepackage[margin=25mm]{geometry}
\usepackage{fontspec}
\setmainfont{Latin Modern Roman}
\usepackage{amsmath,amssymb,mathtools,unicode-math}
\setmathfont{Latin Modern Math}
\usepackage{enumitem,needspace,xurl,hyperref}
\setlength{\parindent}{0pt}
\setlength{\parskip}{5pt}
\setlength{\emergencystretch}{3em}
\newcommand{\N}{\mathbb{N}}
\newcommand{\Z}{\mathbb{Z}}
\newcommand{\Q}{\mathbb{Q}}
\newcommand{\R}{\mathbb{R}}
\newcommand{\C}{\mathbb{C}}
\newcommand{\F}{\mathbb{F}}
\newcommand{\A}{\mathbb{A}}
\newcommand{\PP}{\mathbb{P}}
\newcommand{\E}{\mathbb{E}}
\newcommand{\eps}{\varepsilon}
\newcommand{\dd}{\,\mathrm{d}}
\begin{document}
\section*{UnsolvedMath ID: 30004286}
\textbf{Source problem code:} catalogue entry
\subsection{Definitions and mathematical statement}
Let \(n\) be a fixed positive integer in the degree range intended by the source; that range is not specified in the supplied record. A degree-\(n\) number field is a finite extension \(K/\Q\) with \([K:\Q]=n\). Write \(\mathcal O_K\) for its ring of integers and \(\operatorname{Disc}(K)\) for its signed field discriminant. The field is \emph{monogenic} if there exists \(\theta\in\mathcal O_K\) such that \(\mathcal O_K=\Z[\theta]\). Let \(\widetilde K\) be its normal closure over \(\Q\), and let \(S_n\) denote the symmetric group on \(n\) letters.

Under the conventions of counting fields up to \(\Q\)-isomorphism and bounding the absolute field discriminant, define, for real \(X>0\),
\[
N_n(X)=\#\left\{[K]_{\Q}:\begin{array}{l}
[K:\Q]=n,\quad \mathcal O_K=\Z[\theta]\text{ for some }\theta\in\mathcal O_K,\\
\operatorname{Gal}(\widetilde K/\Q)\cong S_n,\quad |\operatorname{Disc}(K)|<X
\end{array}\right\}.
\]
These counting and discriminant conventions are explicit interpretive choices: the supplied statement does not specify them.

The supplied literature assessment attributes to Bhargava--Shankar--Wang the lower bound
\[
N_n(X)\gg_n X^{1/2+1/n}.
\]
For a fixed admissible degree, this means that there exist constants \(c_n>0\) and \(X_n>0\) such that \(N_n(X)\ge c_nX^{1/2+1/n}\) for every real \(X\ge X_n\). The degree range and any hypotheses of the cited result are not supplied here.

A precise exponent-only interpretation of the proposed optimality is the following conditional formulation: for each fixed degree \(n\) to which the source intends the conjecture to apply, and every real \(\varepsilon>0\), there exist constants \(C_{n,\varepsilon}>0\) and \(X_{n,\varepsilon}>0\) such that
\[
N_n(X)\le C_{n,\varepsilon}X^{1/2+1/n+\varepsilon}
\qquad\text{for every real }X\ge X_{n,\varepsilon}.
\]
Together with the attributed lower bound, this expresses \(N_n(X)=X^{1/2+1/n+o(1)}\) as \(X\to\infty\), with \(n\) fixed. The source does not distinguish this exponent-only assertion from the stronger bound \(N_n(X)\ll_n X^{1/2+1/n}\), or from an asymptotic formula. Consequently, the intended degree range, counting conventions, discriminant convention, and strength of optimality remain unresolved; this is a conditional statement-only draft.
\subsection{Short English statement}
The proposed claim is that the supplied lower-bound exponent \(1/2+1/n\) gives the growth exponent for the number of degree-\(n\) monogenic fields whose normal closures have Galois group \(S_n\). The displayed formulation uses absolute discriminants, counts isomorphism classes, and allows an arbitrarily small positive addition to the upper-bound exponent. These conventions, the intended degrees, and whether a stronger upper bound is sought require source review.
\subsection{Sources}
\begin{itemize}
\item Supplied local UnsolvedMath registry record, permanent ID \(30004286\), title ``Optimal Exponent for Counting Monogenic Number Fields,'' exact source problem code \texttt{catalogue entry}, in the Oberwolfach Reports Open Problems collection. The record lists proposed year 2019 and names no proposer.
\item Supplied report citation: \emph{Analytic Number Theory (2020)}, Oberwolfach Reports, DOI 10.4171/owr/2019/50, pp. 3141--3205. \url{https://doi.org/10.4171/owr/2019/50}. The background describes these reports as documenting mathematical workshops and their open-problem discussions.
\item Bhargava--Shankar--Wang, ``squarefree polynomial discriminants'' (supplied reference label). \url{https://arxiv.org/abs/1611.09806}. The supplied assessment attributes the lower bound and the conjecture of exponent optimality to these authors.
\item Shankar, ``heuristics for \(S_n\)-number fields'' (supplied reference label). \url{https://arxiv.org/abs/2006.09620}. The supplied assessment says later heuristic work supports the expected scale without proving a matching upper bound or asymptotic.
\item Literature assessment dated exactly \texttt{2026-08-22}: the record labels the problem open. This is supplied dated curation derived from the OWR material, not an independently verified assessment; time-sensitive status and cited claims require independent verification.
\item The registry labels the statement \texttt{corrected\_verified}, with the qualification that only formula delimiters, notation, grammar, and typography were repaired and the mathematical content and scope were preserved. This label does not establish mathematical correctness or resolve the formulation issues above.
\end{itemize}
\end{document}
